\section{预层范畴的函子性}\etiquette[5]{I.sec5}

\setcounter{subsection}{-1}
\subsection{}\etiquette[5.0]{I.sec5.0}
设 $C$, $C'$ 和 $D$ 是三个范畴, $u:C\rightarrow C'$ 是一个函子.
由与函子 $u$ 复合所得到的函子记为 $u^\ast$:
\sisi{$$
u^\ast:\Hhom({C'}^\circ,D)\longrightarrow \Hhom(C^\circ, D),
G\longmapsto G\circ u{\quoi},
$$}{$$u^\ast:\left\{\begin{array}{ccc}
                      \Hhom({C'}^\circ,D) & \longrightarrow & \Hhom(C^\circ, D) \\
                      G & \longmapsto & G\circ u{\quoi}
                    \end{array}
\right.$$}
函子 $u^\ast$ 与归纳极限和投射极限交换.

\begin{enonce*}{命题 5.1}\etiquette[5.1]{I.prop5.1}
假设\pageoriginale $C$ 是小范畴, 且 $D$ 中的 $𝒰$-归纳极限
(\resp 投射极限) 可表. 函子 $u^\ast$ 有一个左伴随 $u_{!}$
(\resp 右伴随 $u_\ast$). 因而有同构
\begin{gather*}
\Hom_{\Hhom(C^\circ,D)}(F,u^\ast G)\xrightarrow{\sim}
\Hom_{\Hhom({C'}^\circ,D)}(u_{!}F,G)\\
(\mbox{\resp  } \Hom_{\Hhom(C^\circ,D)}(u^\ast G,F)\simeq
\Hom_{\Hhom({C'}^\circ,D)}(G,u_\ast F){\quoi}).
\end{gather*}
\end{enonce*}

\begin{enonce*}[remark]{证明}
我们只证明左伴随函子的存在性. 命题中 \resp 后的部分便形式地由同构
$$
\Hhom(C^\circ,
D)\xrightarrow{\sim}\Hhom(C,D^\circ)^\circ\xrightarrow{\sim}\Hhom((C^\circ)^\circ,D^\circ)^\circ{\quoi}.
$$
推出.
\end{enonce*}

设 $Y$ 是 $C'$ 的一个对象. 用 $I^Y_u$ 表示如下范畴.
$I^Y_u$ 的对象是二元组 $(X,m)$, 其中 $X$ 是 $C$ 的对象,
$m$ 是一个态射 $Y\rightarrow u(X)$. 设 $(X,m)$ 和 $(X',m')$
是 $I^Y_u$ 的两个对象. 从 $(X,m)$ 到 $(X',m')$ 的态射是满足
$m'=u(\xi)m$ 的态射 $\xi:X\rightarrow X'$. 态射的复合按显然的方式定义.

设 $f:Y\rightarrow Y'$ 是 $C'$ 的一个态射. 由复合, 态射 $f$
定义一个函子 $I^f_u:I^{Y'}_u\rightarrow I^Y_u$.

此外还有一个函子 $\pr_Y:I^Y_u\rightarrow C$, 它将对象 $(X,m)$
对应到对象 $X$. 注意, 此时图
\[\vcenter{
\xymatrix@=1.5cm{
       &I^{Y'}_u \ar[rr]^{I^f_u} \ar[rd]_{\pr_{Y'}}& & I^Y_u \ar[ld]^{\pr_Y}\\
       & &   C   & \\
}}\leqno{(*)}
\]
交换.

现在\pageoriginale 设 $F$ 是 $C$ 上的一个预层, 并置
\hypertarget{sga4.exp.I.equation.5.1.1}{}$$
\displaylines{\rlap{(5.1.1)}\hfill u_{!}
  F(Y)=\varinjlim_{I^Y_u} F\sisi{}{\circ}\pr_Y(.){\quoi}.\hfill}
$$
图 ($^\ast$) 的交换性和归纳极限的函子性使 $u_{!}F$ 成为 $C'$ 上的预层.
下面证明函子 $u_{!}$ 是函子 $u^\ast$ 的一个左伴随. 为此,
我们证明对 $C'$ 上的每个预层 $G$, 都有一个函子同构
$$
\Hom(u_{!}F,G)\xrightarrow{\sim}\Hom(F,u^\ast G){\quoi}.
$$
设 $\xi\in \Hom(u_{!}F,G)$. 那么对 $C$ 的每个对象 $X$, 都有一个态射
$$
\xi_X:u_{!}F(u(X))\longrightarrow G(u(X)){\quoi}.
$$
但 $(u(X),\id_{u(X)})$ 是 $I^{u(X)}_u$ 的一个对象, 而由归纳极限的定义,
有一个典范态射
$$
F(X)\longrightarrow u_{!}F(u(X)){\quoi}.
$$
因此, 对 $C$ 的每个对象 $X$, 我们得到一个态射
$$
\eta_X:F(X)\longrightarrow G(u(X))
$$
它显然关于 $X$ 具有函子性. 从而有一个态射
$$
\eta:F\longrightarrow u^\ast G{\quoi}.
$$

反过来, 设 $\eta\in \Hom(F,u^\ast G)$. 对 $C'$ 的每个对象 $Y$,
由此得到一个函子态射
$$
\eta_Y:F\sisi{}{\circ}\pr_Y\longrightarrow (u^\ast G)\circ\pr_Y{\quoi},
$$
于是, 再与从函子 $(u^\ast G)\circ\pr_Y$ 到常值函子 $G(Y)$ 的显然态射复合,
便得到一个态射
$$
\xi_Y:u_{!}F(Y)\longrightarrow G(Y)
$$
它关于 $Y$ 具有函子性. 从而有一个态射
$$
\xi:u_{!}F\longrightarrow G{\quoi}.
$$

读者\pageoriginale 可验证, 如此定义的两个映射互为逆映射, 证明即告完成.

\begin{enonce*}{命题 5.2}\etiquette[5.2]{I.prop5.2}
假设 $D$ 中的 $𝒰$-归纳极限可表, 有限投射极限可表,
且滤过 $𝒰$-归纳极限与有限投射极限交换. 还假设 $C$ 中的有限投射极限可表,
且函子 $u$ 左正合 $(\hyperref[I.cor2.3.2]{2.3.2})$.
则 $\Hhom(C^\circ,D)$ 和 $\Hhom({C'}^\circ,D)$ 中的有限投射极限可表,
且函子 $u_{!}$ 左正合.
\end{enonce*}

\begin{enonce*}[remark]{证明}
第一个断言显然. 我们证明第二个断言. 由 \hyperref[I.prop5.1]{5.1} 的证明,
对 $C$ 上的每个预层 $F$ 和 $C'$ 的每个对象 $Y$, 都有
$$
u_{!}F(Y)\xrightarrow{\sim} \varinjlim_{I^Y_u}
F\sisi{}{\circ}\pr_Y{\quoi}.
$$
因此只须证明范畴 $(I^Y_u)^\circ$ 满足公理 (PS 1) 和 (PS 2)
\hyperref[I.def2.7]{(2.7)}, 并且这个范畴连通. 验证留给读者.
\end{enonce*}

% END INCLUDED SGA4_English_current_buildable_source_20260906/source/01/drafts/SGA4_I_LINES_1025_1161_English.texfrag


% BEGIN INCLUDED SGA4_English_current_buildable_source_20260906/source/01/drafts/SGA4_I_LINES_1162_1346_English.texfrag
\setcounter{subsection}{2}
\subsection{}\etiquette[5.3]{I.sec5.3}
将上述结果用于 $D$ 是 $𝒰$-集合范畴的情形, 得到由三个函子组成的序列:
$$
u_{!}, u^\ast, u_\ast {\quoi},
$$
这是一个“伴随函子序列”, 其含义是: 对序列中任意两个相邻函子,
右边的函子是另一个的右伴随. 它们的\pageoriginale 主要性质概括在以下命题中:

\begin{enonce*}{命题 5.4}\etiquette[5.4]{I.prop5.4}
设 $C$ 是一个小范畴, $C'$ 是一个 $𝒰$-范畴, $u:C\rightarrow C'$ 是一个函子.
\begin{enumerate}
\renewcommand{\labelenumi}{\theenumi)}
\item 函子 $u^\ast:\widehat{C'}\rightarrow
  \widehat{C}$ 与归纳极限和投射极限交换.

\item 函子 $u_\ast:\widehat{C}\rightarrow
  \widehat{C'}$ 与投射极限交换. 对 $C$ 上的每个预层 $F$ 和 $C'$ 的每个对象 $Y$, 都有:
  $$
  u_\ast F(Y)\xrightarrow{\sim}
  \Hom_{\widehat{C}}(u^\ast(Y),F){\quoi}.
  $$

\hypertarget{sga4.exp.I.proposition.5.4.item.3}{}\item 函子 $u_{!}:\widehat{C}\rightarrow
  \widehat{C'}$ 与归纳极限交换. 函子 $u_{!}$ 仅在相差同构的意义下有定义,
  但总可以选取它使图
\[
\xymatrix@C=2.75cm@R=1.5cm{
  C \ar[r]^{u}\ar[d]_{h}& C' \ar[d]^{h^\prime}\\
C{\widehat{\phantom{\;i}}}\ar[r]^{u_{!}}&C'{\widehat{\phantom{\;i}}}}
\]
($h$ 和 $h'$ 为典范包含函子) 交换.

对 $C$ 上的每个预层 $F$, 都有:
$$
u_{!}F\xrightarrow{\sim} \varinjlim_{C/F} h'\sisi{}{\circ}u{\quoi}.
$$

\item 若 $C$ 中的有限投射极限可表且 $u$ 左正合 $\hyperref[I.cor2.3.2]{(2.3.2)}$,
则函子 $u_{!}$ 左正合.
\end{enumerate}
\end{enonce*}

\begin{enonce*}[remark]{证明}
断言 (1) 显然. 断言 (2) 由 $u_\ast$ 是右伴随函子这一事实,
结合 (\hyperref[I.prop2.11]{2.11}) 和 (\hyperref[I.prop1.4]{1.4}) 得出.
断言 (3) 也是如此, 不过还要应用 (\hyperref[I.prop3.4]{3.4}).
最后, 断言 (4) 就是 \sisi{5.2}({\hyperref[I.prop5.2]{5.2}}),
由 (\hyperref[I.def2.7]{2.7}) 可知它适用于此处.
\end{enonce*}

\begin{enonce*}{命题 5.5}\etiquette[5.5]{I.prop5.5}
设\pageoriginale $C$ 和 $C'$ 是两个小范畴,
$\xymatrix@C=1cm{
 C \ar@<-2.5pt>[r]_{u} & C'\ar@<-2.5pt>[l]_{v}
}$ 是一对函子, 其中 $v$ 是 $u$ 的左伴随. 那么存在与伴随同构相容的同构:
\[
\xymatrix@R=.75cm{
v^{\ast} \ar[rr]^{\sim}&    & u_!      &\\
v_{\ast} \ar[rr]^{\sim}&    & u^{\ast} & . }
\]
\end{enonce*}

\begin{enonce*}[remark]{证明}
只须给出一个同构 $v^\ast\xrightarrow{\sim}
u_{!}$; 另一个同构由伴随关系得到.
设 $F$ 是 $C$ 上的一个预层, $Y$ 是 $C'$ 的一个对象. 则有:
$$
v^\ast F(Y)\xrightarrow{\sim} \Hom(v(Y),F){\quoi}.
$$
其次, 利用 (\hyperref[I.prop3.4]{3.4}):
$$
\Hom(v(Y),F)\xrightarrow{\sim} \varinjlim_{C/F} \Hom(v(Y),.){\quoi}.
$$
但 $v$ 是 $u$ 的左伴随, 因而:
$$
\varinjlim_{C/F}\Hom(v(Y),.)\xrightarrow{\sim}\varinjlim_{C/F}\Hom(Y,u(.)){\quoi}.
$$
再利用 (\hyperref[I.prop5.4]{5.4}.3)), 得到:
$$
\varinjlim_{C/F}\Hom(Y,u(.))\xrightarrow{\sim}
\Hom(Y,u_{!}F)\xrightarrow{\sim} u_{!}F(Y){\quoi}.
$$
因此, 对 $C'$ 的每个对象 $Y$, 我们确定了一个同构
$v^\ast F(Y)\xrightarrow{\sim} u_{!}F(Y)$, 它显然关于 $Y$ 和 $F$
都具有函子性, \cqfd
\end{enonce*}

\begin{enonce*}{推论 5.5.1}\etiquette[5.5.1]{I.cor5.5.1}
设 $u:C\rightarrow C'$ 是一个有左伴随的函子.
函子 $u_{!}:\widehat{C}\rightarrow
\widehat{C'}$ 与投射极限交换
(回顾一下, 由 $\hyperlink{sga4.exp.I.proposition.5.4.item.3}{(1.4.3)}$,
它也与归纳极限交换).
\end{enonce*}

\begin{enonce*}[remark]{注 5.5.2}\etiquette[5.5.2]{I.rem5.5.2}
这样\pageoriginale 便得到一个“\emph{四个}伴随函子的序列”
(\cf \hyperref[I.sec5.3]{5.3}):
$$
v_{!}, v^\ast=u_{!}, v_\ast = u^\ast, u_\ast{\quoi},
$$
因此, 其中的前三个 (\resp 后三个) 与 $\varinjlim$ (\resp $\varprojlim$) 交换.
\end{enonce*}

\begin{enonce*}{命题 5.6}\etiquette[5.6]{I.prop5.6}
假设与 $\hyperref[I.prop5.4]{5.4}$ 相同. 下列条件等价:
\begin{enumerate}
\renewcommand{\theenumi}{\roman{enumi}}
\renewcommand{\labelenumi}{\theenumi)}
\item 函子 $u$ 全忠实.

\item 函子 $u_{!}$ 全忠实.

\item 伴随态射
  $\id_{\widehat{C}}\rightarrow u^\ast u_{!}$ 是一个同构.

\item 函子 $u_\ast$ 全忠实.

\item 伴随态射 $u^\ast u_\ast\rightarrow
  \id_{\widehat{C}}$ 是一个同构.
\end{enumerate}
\end{enonce*}

\begin{enonce*}[remark]{证明}
显然有 ii) $\Leftrightarrow$ iii) 和 iv) $\Leftrightarrow$
v) (伴随函子的一般性质), 以及 ii)
$\Rightarrow$ i) (\hyperref[I.prop5.4]{5.4}.3)). 我们证明 i) $\Rightarrow$
iii). 函子 $\id_{\widehat{C}}$, $u^\ast$ 和
$u_{!}$ 都与归纳极限交换. 因此由 \hyperref[I.prop3.4]{(3.4)},
只须在 $H$ 可表时证明 $H\rightarrow u^\ast u_{!} H$ 是同构, 这是显然的.
下面证明 iii) 与 v) 等价.
对 $\widehat{C}$ 的每个对象 $H$ (\resp $K$), 用
$\Phi(H):H\rightarrow u^\ast u_{!} H$ (\resp
$\Psi(K):u^\ast u_\ast\rightarrow K$) 表示伴随态射. 我们便有交换图:
\[
\xymatrix@C=3.5cm@L=6pt{
\Hom_{\widehat{C}}(H,u^{\ast}u_{\ast}
K)\ar@{-}[d]_{\wr}\ar[r]^{\quad \Hom(H,\Psi(K))}   & \Hom_{\widehat{C}}(H,K)\\
\Hom_{\widehat{C}}(u_{!}H,u_{\ast}K)\ar@{-}[d]_{\wr} &  \\
\Hom_{\widehat{C}}(u^{\ast}u_{!}H,K)\ar[uru]_{\quad
   \Hom(\Phi(H),K)} & .
}
\]

由此\pageoriginale 可见, 对每个 $H$, $\Phi(H)$ 都是同构,
当且仅当对每个 $K$, $\Psi(𝒦)$ 都是同构, \cqfd
\end{enonce*}

\begin{enonce*}[remark]{注 5.7}\etiquette[5.7]{I.rem5.7}
\begin{enumerate}
\renewcommand{\theenumi}{\alph{enumi}}
\renewcommand{\labelenumi}{\theenumi)}
\item 等价关系 ii) $\Leftrightarrow$ iii)
  $\Leftrightarrow$ iv) $\Leftrightarrow$ v) 是伴随函子的一般结论.

\item \hyperref[I.prop1.4]{1.4} 的证明中给出的 $u_{!}$ \resp $u_\ast$
的显式形式直接表明, 在 \hyperref[I.def1.1]{1.1} 的一般假设下,
若 $u$ 全忠实, 则伴随态射 $\id\rightarrow u^\ast u_{!}$
(\resp $u^\ast u_\ast\rightarrow \id$) 是同构,
\ie $u_{!}$ (\resp $u_\ast$) 全忠实.
\end{enumerate}
\end{enonce*}

\setcounter{subsection}{8} \setcounter{subsubsection}{-1}

% END INCLUDED SGA4_English_current_buildable_source_20260906/source/01/drafts/SGA4_I_LINES_1162_1346_English.texfrag


% BEGIN INCLUDED SGA4_English_current_buildable_source_20260906/source/01/drafts/SGA4_I_LINES_1347_1496_English.texfrag
\subsubsection{}\etiquette[5.8.0]{I.sec5.8.0}
设 $\gamma$ 是一个由有限投射极限定义的代数结构类型,
$𝒰$-$\gamma$-$\Ens$ 是 $𝒰$-$\Ens$ 中的 $\gamma$-对象范畴,
esj: $𝒰$-$\gamma$-$\Ens\rightarrow 𝒰$-$\Ens$ 是底层集合函子
(为简单起见, 我们假设所考虑的结构类型只有一个底层集合).
设 $C$ 是一个范畴. 与 $\esj$ 复合得到一个函子, 记为
$$
\mbox{esj}^{\widehat{\phantom{,}}}:\Hhom(C^\circ,
𝒰\text{-}\gamma\text{-}\Ens)\longrightarrow
\widehat{C}{\quoi}.
$$
由于 $\widehat{C}$ 中的投射极限逐点计算,
函子 $\mbox{esj}^{\widehat{\phantom{,}}}$ 分解为一个等价
\hypertarget{sga4.exp.I.equation.5.8.1}{}$$
\displaylines{\rlap{(5.8.1)}\hfill \Hhom(C^\circ,
  𝒰\text{-}\gamma\text{-}\Ens)\xrightarrow{\approx}
  \widehat{C}_\gamma{\quoi},\hfill}
$$
其中 $\widehat{C}_\gamma$ 是 $\widehat{C}$ 中的 $\gamma$-对象范畴,
再接上一个仍记为
$$
\mbox{esj}^{\widehat{\phantom{,}}}:\widehat{C}_\gamma\longrightarrow
\widehat{C}{\quoi},
$$
的函子, 称为“\emph{底层集合预层}”函子.

\setcounter{subsubsection}{1}
\subsubsection{}\etiquette[5.8.2]{I.sec5.8.2}
假设\pageoriginale 函子
$\esj:𝒰$-$\gamma$-$\Ens\rightarrow 𝒰$-$\Ens$ 有一个左伴随
$\Lib:𝒰$-$\Ens\rightarrow 𝒰$-$\gamma$-$\Ens$
\footnote[1]{函子 $\Lib$ 是“自由生成的 $\gamma$-对象”函子.
例子: 自由群, 自由 Abel 群, 自由 $A$-模, \etc}
(事实上可以证明, 这个条件总是满足的). 与 $\Lib$ 复合得到一个函子
$$
\Lib^{\widehat{\phantom{,}}}:\widehat{C}\longrightarrow
\Hhom(C^\circ, 𝒰\text{-}\gamma\text{-}\Ens)
$$
再与等价 \hyperlink{sga4.exp.I.equation.5.8.1}{(5.8.1)} 复合,
得到一个仍记为
$$
\Lib^{\widehat{\phantom{,}}}:\widehat{C}\longrightarrow
\widehat{C}_\gamma
$$
的函子, 称为“自由生成的 $\gamma$-对象预层”函子.
函子 $\Lib^{\widehat{\phantom{,}}}$ 是函子
$\esj^{\widehat{\phantom{,}}}$ 的左伴随.

\begin{enonce*}{命题 5.8.3}\etiquette[5.8.3]{I.prop5.8.3}
设 $\gamma$ 是由有限投射极限定义的代数结构类型,
并且 $𝒰$-$\Ens$ 中的 $\gamma$-对象范畴的 $𝒰$-归纳极限可表
\footnote[2]{可以证明, 这个条件总是满足的.};
设 $C$ 是属于 $𝒰$ 的一个范畴, $C'$ 是一个 $𝒰$-范畴,
$u:C\rightarrow C'$ 是一个函子. 用 $\widehat{C}_\gamma$
(\resp $\widehat{C'}_\gamma$) 表示 $\widehat{C}$
(\resp $\widehat{C'}$) 中的 $\gamma$-对象范畴,
用 ${u^\ast}^\gamma$ 表示函子 $u^\ast$ 在 $\gamma$-对象上诱导的函子.
由 \hyperref[I.prop5.1]{5.1} 和等价
\hyperlink{sga4.exp.I.equation.5.8.1}{5.8.1} 可知,
函子 ${u^\ast}^\gamma$ 有一个左 (\resp 右) 伴随.
这个函子记为 $u_{!_\gamma}$ (\resp $u_{\ast_\gamma}$).
\begin{enumerate}
\item 函子 ${u^\ast}^\gamma$ 与归纳极限和投射极限交换. 图
\[\vcenter{
\xymatrix@C=3cm@R=1.5cm{
 \widehat{C'}_\gamma \ar[r]^{u^{\ast^{\gamma}}} \ar[d]_{\esj'\widehat{\phantom{a}}}&
\widehat{C}_\gamma  \ar[d]_{\esj^{\widehat{\phantom{a}}}}\\
 \widehat{C'} \ar[r]^{u_{\ast}}& \widehat{C}
}}\leqno{(*)}
\]
\pageoriginale 交换 ($\esj'^{\widehat{\phantom{,}}}$ 和
$\esj^{\widehat{\phantom{,}}}$ 表示“底层集合”函子).

\item 函子 $u_{\ast_\gamma}$ 与投射极限交换. 图
\[\vcenter{
\xymatrix@C=3cm@R=1.5cm{
 \widehat{C}_\gamma \ar[r]^{u_{\ast_\gamma}} \ar[d]_{\esj\widehat{\phantom{a}}}&
  \widehat{C'}_\gamma\ar[d]_{\esj'\widehat{\phantom{a}}}\\
\widehat{C} \ar[r]^{u_{\ast}}&
\widehat{C'} }}\leqno{({*}{*})}
\]
在相差同构的意义下交换.

\item 函子 $u_{!_\gamma}$ 与归纳极限交换. 假设
$\esj^{\widehat{\phantom{,}}}$
(\resp ${\esj'}^{\widehat{\phantom{,}}}$) 有一个左伴随
$\Lib^{\widehat{\phantom{,}}}$
(\resp ${\Lib'}^{\widehat{\phantom{,}}}$) $\hyperref[I.sec5.8.2]{(5.8.2)}$.
图
\[\vcenter{
\xymatrix@C=3.5cm@R=1.75cm{
 \widehat{C} \ar[r]^{u_{!}} \ar[d]_{\Lib\widehat{\phantom{a}}}&
  \widehat{C'} \ar[d]_{\Lib'\widehat{\phantom{a}}}\\
\widehat{C}_{\gamma} \ar[r]^{u_{!\gamma}} &
\widehat{C'}_{\gamma} }}\leqno{({*}{*}{*})}
\]
在相差同构的意义下交换.
\end{enumerate}

假设 $u_{!}$ 与有限投射极限交换
(\hyperref[I.prop5.4]{5.4}) 和 (\hyperref[I.prop5.6]{5.6}). 则图
\[\vcenter{
\xymatrix@C=3.5cm@R=1.5cm{
 \widehat{C}_\gamma \ar[r]^{u_{!\gamma}} \ar[d]_{\esj\widehat{\phantom{a}}}&
  \widehat{C'}_\gamma\ar[d]_{\esj'\widehat{\phantom{a}}}\\
\widehat{C} \ar[r]^{u_{!}}& \widehat{C'}
}}\leqno{({*}{*}{*}{*})}
\]
在相差同构的意义下交换, 且 $u_{!\gamma}$ 与有限投射极限交换.
\end{enonce*}

\begin{enonce*}[remark]{证明}
断言 (1) 显然. 断言 (2) 也显然, 因为 $u_\ast$ 是右伴随,
所以 (\hyperref[I.prop2.11]{2.11}) 与投射极限交换, 尤其与有限投射极限交换;
这就给出了图 $({\ast}{\ast})$ 的交换性.
图 $({\ast}{\ast}{\ast})$ 的交换性来自左伴随函子在\pageoriginale 相差同构意义下的唯一性;
最后, 图 $({\ast}{\ast}{\ast}{\ast})$ 的交换性直接来自 $u_{!}$ 与有限投射极限交换这一事实.
\end{enonce*}

\begin{enonce*}[remark]{记号 5.9}\etiquette[5.9]{I.not5.9}
我们常滥用记号, 将函子 ${u^\ast}^\gamma$ 和 $u_{\ast\gamma}$
记为 $u^\ast$ 和 $u_\ast$. 由于图 $(\ast)$ 和 $({\ast}{\ast})$ 交换,
这不会引起混淆. 另一方面, 当 $u_{!}$ 不与有限投射极限交换时,
图 $({\ast}{\ast}{\ast}{\ast})$ 在相差同构的意义下也不交换,
因而必须使用记号 $u_{!}$ 和 $u_{!\gamma}$ 来避免混淆.
\end{enonce*}

\setcounter{subsection}{9}

% END INCLUDED SGA4_English_current_buildable_source_20260906/source/01/drafts/SGA4_I_LINES_1347_1496_English.texfrag


% BEGIN INCLUDED SGA4_English_current_buildable_source_20260906/source/01/drafts/SGA4_I_LINES_1497_1738_English.texfrag
\subsection{}\etiquette[5.10]{I.sec5.10}
设 $C$ 是一个小范畴. 对 $\widehat{C}$ 的每个对象 $X$,
用 $C/X$ 表示以 $X$ 为靶、以 $C$ 的对象为源的箭头组成的范畴.
源函子定义一个函子 $j_X:C/X\rightarrow C$.
设 $m:Y\rightarrow X$ 是 $\widehat{C}$ 的一个态射.
态射的复合定义一个函子 $j_m:C/Y\rightarrow C/X$. 图
\[
\xymatrix@=1.5cm{
C/Y \ar[rr]^{j_m} \ar[rd]_{j_{Y}}&   & C/X\ar[ld]^{j_{X}}\\
    & C &
}
\]
交换. 由 \hyperref[I.prop5.1]{5.1} 可知, 对 $(C/X)^{\widehat{\phantom{,}}}$
的每个对象 $F$ 和 $C$ 的每个对象 $Y$, 都有
\hypertarget{sga4.exp.I.equation.5.10.1}{}$$
\displaylines{\rlap{(5.10.1)}\hfill j_{X!}F(Y)=\coprod_{u\in
    \Hom_{\widehat{C}}(Y,X)} F(u){\quoi}.\hfill}
$$
当 $C$ 是一个 $𝒰$-范畴时, 公式
\hyperlink{sga4.exp.I.equation.5.10.1}{(5.10.1)} 定义 $j_{X!}$,
而且可验证如此定义的函子 $j_{X!}$ 总是函子
$j^\ast_X:\widehat{C}\rightarrow
(C/X)^{\widehat{\phantom{,}}}$ (\hyperref[I.sec5.0]{5.0}) 的右伴随\pageoriginale.

\begin{enonce*}{命题 5.11}\etiquette[5.11]{I.prop5.11}
设 $C$ 是一个 $𝒰$-范畴, $X$ 是 $C$ 上的一个预层.
\begin{enumerate}
\renewcommand{\labelenumi}{\theenumi)}
\item 函子
$$
j_{X!}:(C/X)^{\widehat{\phantom{,}}}\longrightarrow
\widehat{C}
$$
经过范畴 $\widehat{C}/X$ 分解:
$$
(C/X)^{\widehat{\phantom{,}}}\xrightarrow{e_X}
\widehat{C}/X\longrightarrow
\widehat{C}{\quoi}.
$$
函子 $e_X$ 是一个范畴等价.

\item 函子 $e_{X}\circ j^\ast_X:\widehat{C}
\to  \widehat{C}/X$ 与函子
$H\mapsto (H\times X \xrightarrow{\pr_2} X)$ 典范同构.
\end{enumerate}
\end{enonce*}

\begin{enonce*}[remark]{证明}
\begin{enumerate}
\renewcommand{\labelenumi}{\theenumi)}
\item 设 $f$ 是 $(C/X)^{\widehat{\phantom{,}}}$ 的终对象.
有一个典范同构 $f\xrightarrow{\sim}
\varinjlim_{Y\in \ob C/X}Y$, 因而
$j_{X!}(f)=\varinjlim_{Y\in \ob C/X}j_X(Y)$ \hyperref[I.prop5.4]{(5.4)}.
但 $j_{X!}(f)\simeq X$, 这就给出了所述分解.
要证明 $e_X$ 是等价, 只须给出一个拟逆函子.
将 $\widehat{C}/X$ 的每个对象 $H\rightarrow X$ 对应到 $C/X$ 上的预层
$$
(Y\rightarrow X)\longmapsto
  \Hom_{\widehat{C}/X}((Y\rightarrow X), (H\rightarrow
  X)){\quoi}.
$$

\item 函子 $e_X\circ j^\ast_X$ 是遗忘函子的右伴随,
因此函子 $e_X\circ j^\ast$ 与函子 $H\mapsto (H\times
  H\xrightarrow{\pr_2}X)$ 典范同构.
\end{enumerate}
\end{enonce*}

\setcounter{subsection}{11}
\subsection{}\etiquette[5.12]{I.sec5.12}
设 $m:Y\rightarrow X$ 是 $\widehat{C}$ 的一个态射.
由 \hyperref[I.prop5.11]{(5.11)}, 态射 $m$ 典范同构于
$(C/X)^{\widehat{\phantom{,}}}$ 的某个对象在 $e_X$ 下的像,
我们将这个对象记为 $[m]$. 将函子 $e_X$ 限制到相应子范畴上, 得到一个等价
$$
(C/X)/[m]\xrightarrow{e_m}C/Y{\quoi}.
$$
图\pageoriginale
\[
\xymatrix{
  (C/X)/[m]\ar[rr]^{e_m} \ar[rd]_{j_{[m]}} &       & C/Y \ar[ld]^{j_m}\\
          & C/X   &
}
\]
在相差典范同构的意义下交换.

\subsection{}\etiquette[5.13]{I.sec5.13}

我们记录一个将在第 VI 讲中用到的结果.
设 $u:C\rightarrow C'$ 是小范畴之间的一个函子.
对 $\widehat{C}$ 的每个对象 $H$, 用 $u/H:C/H\rightarrow C'/u_{!}H$
表示将每个态射 $m:X\rightarrow H$ 对应到态射
$u_{!}X\xrightarrow{u_{!}m}u_{!}H$ 的函子
(我们知道 (\hyperref[I.prop5.4]{5.4}), 总可以令 $u_{!}X=uX$).
下图交换:
\hypertarget{sga4.exp.I.equation.5.13.1}{}\[\vcenter{
\xymatrix@C=3cm@R=1.75cm{
C/H \ar[r]^{u/H}\ar[d]_{J_H} & C'/u_{!}H \ar[d]_{j_{u_{!}H}} \\
C\ar[r]^{u}   & C'\quad. }}\leqno{(5.13.1)}
\]
因此有一个在相差同构意义下交换的图:
%\[\vcenter{
%\xymatrix@C=2.5cm@R=1.5cm{ (C/H){\widehat{\phantom{a}}}
%\ar[r]^{(u/H)_{!}} \ar[d]_{(j_H)_!}&
%  (C'/u^·H){\widehat{\phantom{a}}} \ar[d]_{(j_{u!H})_{!}}\\
%  C{\widehat{\phantom{\;a}}} \ar[r]^{u_!}& C'{\widehat{\phantom{\;a}}}\quad,
%}}\leqno{(5.13.2)}
%\]
又因为 $(u/H)_{!}$ 将 $(C/H)^{\widehat{\phantom{,}}}$ 的终对象映到
$(C'/u_{!}H)^{\widehat{\phantom{,}}}$ 的终对象, 所以图
\hypertarget{sga4.exp.I.equation.5.13.3}{}\[\vcenter{
\xymatrix@C=3.5cm@R=1.5cm{ (C/H)\widehat{\phantom{a}}
\ar[r]^{(u/H)_!} \ar[d]_{e_H}&
  (C'/u_!H)\widehat{\phantom{a}} \ar[d]_{e_{u_{!}H}}\\
C\widehat{\phantom{a}}/H  \ar[r]^{u_{!}/H} &
C'\widehat{\phantom{a}}/u_{!}H }}\leqno{(5.13.3)}
\]
在相差同构的意义下\pageoriginale 交换 \hyperref[I.def6.1]{(6.1)}.

\begin{enonce*}{命题 5.14}\etiquette[5.14]{I.prop5.14}
设 $u:C\rightarrow C'$ 是小范畴之间的一个函子. 假设 $u$ 具有如下性质:

{\rm(PPF)} 对 $C$ 的每个对象 $X$, 函子
$u/X:C/X\rightarrow C'/uX$ 全忠实.

则:
\begin{enumerate}
\renewcommand{\labelenumi}{\theenumi)}
\item 设 $f$ 是 $\widehat{C}$ 的终对象. 函子 $u$ 分解为
$$
C=C/f\xrightarrow{u/f}C'/u_{!}f\xrightarrow{j_{u_{!}f}} C'{\quoi}.
$$
函子 $u/f$ 全忠实.

\item 函子 $u_{!}:\widehat{C}\rightarrow
  \widehat{C'}$ 分解为
$$
\widehat{C}=(C/f)^{\widehat{\phantom{,}}}\xrightarrow{(u/f)_{!}}
(C'/u_{!}f)^{\widehat{\phantom{,}}}\xrightarrow{e_{u_{!}f}}\widehat{C'}/u_{!}f\longrightarrow
\widehat{C'}{\quoi},
$$
其中函子 $(u/f)_{!}$ 全忠实, 函子 $e_{u_{!}f}$ 是一个等价,
函子 $\widehat{C'}/u_{!}f\rightarrow\widehat{C'}$ 是遗忘函子.

\item 特别地, 函子 $u_{!}$ 忠实, 因而单位态射
$\id\xrightarrow{\Phi}u^\ast u_{!}$ 是单态射.
此外, 对 $\widehat{C}$ 的每个态射 $\alpha:H\rightarrow K$, 图
\[
\xymatrix@C=3.5cm@R=1.5cm@M=.2cm{
H \ar@{^{(}->}[r]^{\Phi(H)}\ar[d]_{\alpha} & u^{*}u_{!}H \ar[d]^{u.u^{\bigdot}(\alpha)}\\
K \ar@{^{(}->}[r]^{\Phi(K)} &  u^{*}u_!K }
\]
是 Cartesian 图.
\end{enumerate}
\end{enonce*}

\begin{enonce*}[remark]{证明}
\begin{enumerate}
\renewcommand{\labelenumi}{\theenumi)}
\item 所述\pageoriginale 分解来自图
\hyperlink{sga4.exp.I.equation.5.13.1}{(5.13.1)}.
函子 $u$ 忠实. 因而 $u/f$ 忠实. 我们证明它全忠实.
设 $X$ 和 $Y$ 是 $C$ 的两个对象,
$\can_X:uX\rightarrow u_{!}f$ (\resp $\can_Y:uY\rightarrow u_{!}f$)
是典范态射, 并设
\[
\xymatrix@=1.25cm{
uX \ar[rr]^{m} \ar[rd]_{\can_X}&               & uY \ar[ld]^{\can_Y}\\
   & u^{\bigdot}f &
}
\]
是 $C'/u_{!}f$ 的一个态射. 我们有
$u_{!}f=\varinjlim_{Z\in \ob C}uZ$, 因而 \hyperref[I.prop3.1]{(3.1)}
$$
\Hom_{\widehat{C'}}(uX,u_{!}f)=\varinjlim_{Z\in
  \ob C}\Hom_{C'}(uX, uZ){\quoi}.
$$
由归纳极限的定义, 说 $\can_Ym=\can_X$, 等价于说存在
\begin{enumerate}
\renewcommand{\theenumii}{\alph{enumii}}
\renewcommand{\labelenumii}{\theenumii)}
\item $C$ 的一个有限对象序列 $X_i$, $i\in [0,n]$,
满足 $X_\circ=X$, $X_n=Y$,

\item 对每个 $i$, 一个态射 $m_i:uX\rightarrow
  uX_i$ $(m_\circ=\id_{X}, m_n=m)$,

\item 对每一对 $(i,i+1)$, 一个态射 $f_i:X_i\rightarrow
  X_{i+1}$ 或 $f_i:X_{i+1}\rightarrow X_i$,
\end{enumerate}
使图
\begin{displaymath}
\xymatrix@R=1.25cm@C=1.25cm{
   &  uX \ar[rd]^{\;\;m_{i+1} \qquad\quad\text{ 或 }} \ar[ld]_{m_i}&\\
  uX_i\ar[rr]^{u(f_i)} &          &  uX_{i+1}
}\qquad \xymatrix@R=1.25cm@C=1.25cm{
      &  uX \ar[rd]^{\;\;m_{i+1}} \ar[ld]_{m_i}&\\
 uX_i &                                        &
   uX_{i+1}\ar[ll]_{u(f_i)}\quad ,
}
\end{displaymath}
交换.

于是对 $i$ 归纳并利用性质 (PPF), 立即得出 $m_i$ 具有形式 $u(p_i)$.
特别地, $m=u(p)$, 因而 $u/f$ 全忠实.

\item 所述\pageoriginale 分解显然. 由 \hyperref[I.prop5.6]{5.6},
函子 $(u/f)_{!}$ 全忠实. 其他断言由 \hyperref[I.prop5.11]{5.11} 得出.

\item 函子 $u_{!}$ 是忠实的遗忘函子与全忠实函子的复合, 因而忠实.
由伴随函子的一般性质可知, 单位态射 $\id\rightarrow u^\ast u_{!}$ 是单态射.
由 2), 函子 $u_{!}$ 可看成一个全忠实函子
$v:\widehat{C}\to\widehat{C'}/u_{!}f$ 与遗忘函子的复合.
因此, $u_{!}$ 的右伴随函子 $u^\ast$ 是“与 $u_{!}f$ 作积”函子
(遗忘函子的右伴随) 与 $v$ 的右伴随函子 $w$ 的复合.
此外, 由于 $v$ 全忠实, 单位态射 $\id\rightarrow wv$ 是同构.
最后一个断言容易得出.
\end{enumerate}
\end{enonce*}


% END INCLUDED SGA4_English_current_buildable_source_20260906/source/01/drafts/SGA4_I_LINES_1497_1738_English.texfrag


% BEGIN INCLUDED SGA4_English_current_buildable_source_20260906/source/01/drafts/SGA4_I_LINES_1739_1946_English.texfrag
\section{忠实函子与保守函子}\etiquette[6]{I.sec6}\pageoriginale

\begin{enonce*}{定义 6.1}\etiquette[6.1]{I.def6.1}
设 $E$ 是一个范畴, $(\varphi_i:E\rightarrow F_i)_{i\in I}$ 是一个函子族
$$
\varphi_i:E\longrightarrow F_i{\quoi}.
$$
若对 $E$ 的每对对象 $X$, $Y$ 以及每对箭头
$u$, $v:X\rightrightarrows Y$, 对每个 $i\in I$ 都有
$\varphi_i(u)=\varphi_i(v)$ 蕴含 $u=v$
(换言之, 由 $(\varphi_i)$ 定义的映射 $\Hom(X,Y)\rightarrow
\Pi_i\Hom(\varphi_i(X),\varphi_i(Y))$ 是单射),
就称函子族 $(\varphi_i)$ \emph{忠实}.
若 $E$ 的每个满足对所有 $i\in I$, $\varphi_i(u)$ 都是同构的箭头 $u$
本身也是同构, 就称函子族 $(\varphi_i)$ \emph{保守}.
若只要 $u$ 是单态射 (\resp 满态射, \resp \ldots), 上述条件就成立,
则称 $(\varphi_i)$ 对单态射 (\resp 满态射, \resp \ldots) 保守.
\end{enonce*}

\setcounter{subsection}{1} \setcounter{subsubsection}{0}
\subsubsection{}\etiquette[6.1.1]{I.sec6.1.1}
若引入由函子族 $(\varphi_i)$ 定义的唯一函子
$$
\varphi:E\longrightarrow F=\prod_{i\in I} F_i
$$
则显然, 这个族忠实 (\resp 保守, \resp 对单态射保守, \resp\ldots),
当且仅当函子 $\varphi$ 忠实 (\resp 保守, \resp\ldots)
(这里的意思是仅由 $\varphi$ 组成的族忠实, \resp 保守, \resp\ldots).
因此, 下文即使只考虑由单个函子组成的族, 也没有多大不便.
不过, 为便于以后引用, 以下结果仍对函子族来陈述.

当\pageoriginale $\varphi_i$ 满足适当的正合性条件时,
\hyperref[I.def6.1]{6.1} 中的概念尤其有用, 而且在这种情况下它们往往重合:

\begin{enonce*}{命题 6.2}\etiquette[6.2]{I.prop6.2}
记号与 \hyperref[I.prop5.1]{5.1} 相同.
\begin{enumerate}
\renewcommand{\theenumi}{\roman{enumi}}
\renewcommand{\labelenumi}{(\theenumi)}
\item 若 $E$ 中双箭头的核或双箭头的余核可表,
且 $\varphi_i$ 与它们交换, 则有蕴含关系
$$
(\varphi_i)\mbox{ 保守 } \Longrightarrow (\varphi_i)\mbox{
    忠实}{\quoi}.
$$

\item 假设 $E$ 中的纤维积 (\resp 纤维余积) 可表,
且 $\varphi_i$ 与它们交换. 假设 $(\varphi_i)$ 忠实或保守;
则对 $E$ 的每个箭头 $u$, $u$ 是单态射 (\resp 满态射),
当且仅当对每个 $i\in I$, $\varphi_i(u)$ 都如此.

\item 假设 $E$ 中的纤维积和纤维余积可表, $\varphi_i$ 与它们交换,
且 $E$ 中每个既是单态射又是满态射的箭头 (即双态射) 都是同构.
则有蕴含关系
$$
(\varphi_i)\mbox{ 忠实 } \Longrightarrow (\varphi_i)\mbox{
    保守}{\quoi}.
$$

\item 假设 $E$ 中的纤维积 (\resp 纤维余积) 可表,
且 $\varphi_i$ 与它们交换. 则若 $(\varphi_i)$ 对单态射
(\resp 满态射) 保守, 它实际上就是保守的.

\item 设 $D$ 是一种图式类型, $F:d\mapsto F(d)$ 是 $E$ 中一个 $D$ 型图式,
$X$ 是 $E$ 的一个对象, $u=(u_d)_{d\in D}$ 是一个箭头族
$$
X\longrightarrow F(d)\quad (\mbox{\resp  } F(d)\longrightarrow
X){\quoi}.
$$
假设\pageoriginale $(\varphi_i)$ 保守, $E$ 中 $D$ 型的投射极限
(\resp 归纳极限) 可表, 且 $\varphi_i$ 与它们交换.
则 $u$ 使 $X$ 成为 $F$ 在 $E$ 中的投射极限 (\resp 归纳极限),
当且仅当对每个 $i\in I$, $\varphi_i(u)$ 使 $\varphi_i(X)$ 成为
$\varphi_i(D)$ 在 $E_i$ 中的投射极限 (\resp 归纳极限).
\end{enumerate}
\end{enonce*}

\noindent \emph{证明.}
\begin{enumerate}
\renewcommand{\theenumi}{\roman{enumi}}
\renewcommand{\labelenumi}{(\theenumi)}
\item 对于不带 \resp 的陈述, 给定一个双箭头
$u$, $v:X\rightrightarrows Y$, 只须将等式 $u=v$ 表述为
包含态射 $\Ker(u,v)\rightarrow X$ 是同构这一条件.
此处及下文中, 对偶陈述的对偶论证均不再重复.

\item 若 $(\varphi_i)$ 忠实, 则 $u:X\rightarrow Y$ 是单态射这一条件,
可表述为纤维积 $\fprod{X}{X}{Y}$ 的两个投射满足等式 $\pr_1=\pr_2$.
若 $(\varphi_i)$ 保守, 则可表述为对角态射
$\delta:X\rightarrow \fprod{X}{X}{Y}$ 是同构这一条件.

\setcounter{enumi}{3}
\item 由于后一论证中的态射 $\delta$ 是单态射,
可见实际上只须假设 $(\varphi_i)$ 对单态射保守.
但这又蕴含 $(\varphi_i)$ 无条件地保守.
事实上, 若 $u\in F\ell\,E$ 满足 $\varphi_i(u)$ 都是同构,
则前面的论证表明它们是单态射, 从而由关于 $(\varphi_i)$ 的假设,
它们是同构.

\setcounter{enumi}{2}
\item 这是 (ii) 的显然推论.

\setcounter{enumi}{4}
\item 这是定义的显然推论.
\end{enumerate}

我们记录由 (i), (ii) 和 (iv) 得到的以下推论:

\begin{enonce*}{推论 6.3}\etiquette[6.3]{I.cor6.3}
假设\pageoriginale $E$ 中的纤维积和纤维余积可表,
$\varphi_i$ 与它们交换, 并且双箭头的核或双箭头的余核可表,
$\varphi_i$ 与它们交换.
(例如, $E$ 中的有限投射极限和有限归纳极限可表,
且 $\varphi_i$ 是正合函子, 就足够了.) 则下列条件等价:
\begin{itemize}
\item[a)] $(\varphi_i)$ 忠实.

\item[b)] $(\varphi_i)$ 保守.

\item[c)] $(\varphi_i)$ 对单态射保守.

\item[${\rm c')}$] $(\varphi_i)$ 对满态射保守.
\end{itemize}
\end{enonce*}
为便于引用, 我们还记录:

\begin{enonce*}{命题 6.4}\etiquette[6.4]{I.prop6.4}
设 $\varphi:E\rightarrow F$ 是一个有右伴随 $\psi$ 的函子
(因而 $\Hom(\varphi(X),Y)\simeq \Hom(X,\psi(Y))$).
函子 $\varphi$ (\resp $\psi$) 忠实,
当且仅当对 $E$ 的每个元素 $X$ (\resp $F$ 的每个元素 $Y$),
单位态射 $X\rightarrow \psi\varphi(X)$ 是单态射.
函子 $\varphi$ (\resp $\psi$) 全忠实,
当且仅当上述单位态射是同构.
\end{enonce*}
事实上, 若 $X'$, $X$ 是 $E$ 的两个对象, 则映射
\begin{equation*}
\Hom(X',X)\longrightarrow
  \Hom(\varphi(X'),\varphi(X))\tag{$\ast$}
\end{equation*}
可视为对
\begin{equation*}
X\longrightarrow \psi
  \varphi (X)\tag{${\ast}{\ast}$}
\end{equation*}
应用函子 $\Hom(X',-)$ 所得的映射.
因此, 固定 $X$ 后, 对每个 $X'$ 映射 $(\ast)$ 都是单态射\pageoriginale
(\resp 同构) 的充要条件, 是单位态射 $({\ast}{\ast})$ 为单态射
(\resp 同构).

\begin{enonce*}{命题 6.5}\etiquette[6.5]{I.prop6.5}
设 $p:E'\rightarrow E$ 是一个函子. 下列条件等价:
\begin{enumerate}
\renewcommand{\theenumi}{\roman{enumi}}
\renewcommand{\labelenumi}{(\theenumi)}
\item $p$ 忠实、保守且纤维化 {\rm(SGA 1 VI 6.1).}

\item $p$ 是一个纤维为离散范畴的纤维化函子.

\item 对每个 $X'\in \ob E'$, 由 $p$ 诱导的函子 $E'_{/X'}\rightarrow
  E_{/p(X)}$ 是一个在对象上满射的范畴等价.

\item (当 $E$ 是一个 $𝒰$-范畴时.) 存在元素
  $F\in \ob\widehat{C}$ 和 $E$ 上的一个范畴等价
  {\rm(SGA 1 VI 4.3)} $E'\xrightarrow{\approx}E_{/F}$
  (其中 $E_{/F}$ 是 $\widehat{E}_{/F}$ 的全子范畴,
  由源属于 $E$ 的箭头 $X\rightarrow F$ 组成).
\end{enumerate}
\end{enonce*}

\setcounter{subsection}{5} \setcounter{subsubsection}{0}

% END INCLUDED SGA4_English_current_buildable_source_20260906/source/01/drafts/SGA4_I_LINES_1739_1946_English.texfrag


% BEGIN INCLUDED SGA4_English_current_buildable_source_20260906/source/01/drafts/SGA4_I_LINES_1947_2034_English.texfrag
\subsubsection{}\etiquette[6.5.1]{I.sec6.5.1}
回顾一下, 称一个范畴 $C$ 是\emph{离散的}, 是指它是一个\emph{群胚}
(\ie 其中的每个箭头都可逆), 并且是\emph{刚性的}
(\ie 每个对象的自同构群都退化为平凡群);
等价地, 这个范畴等价于某个集合 $I$ 定义的范畴 $C'$
(其中 $\text{ob}C'=I$, 且仅有恒等箭头).
若已假设 $C$ 是群胚, 则说 $C$ 离散, 等价于说对 $C$ 的两个对象 $X$, $Y$,
从 $X$ 到 $Y$ 至多存在一个箭头, \ie $C$ 同构于一个预序集所定义的范畴.

\hyperref[I.prop6.5]{6.5} 中条件 (i) 和 (ii) 的等价性\pageoriginale
是上述观察以及以下引理的直接推论:

\begin{enonce*}{引理 6.5.2}\etiquette[6.5.2]{I.lem6.5.2}
设 $p:E'\rightarrow E$ 是一个纤维化函子. 则:
\begin{enumerate}
\renewcommand{\theenumi}{\roman{enumi}}
\renewcommand{\labelenumi}{(\theenumi)}
\item $p$ 保守的充要条件是它的纤维范畴都是群胚.

\item $p$ 忠实的充要条件是它的纤维范畴都是有序范畴.
\end{enumerate}
\end{enonce*}

\noindent \emph{\hyperref[I.lem6.5.2]{6.5.2} 的证明.}
\begin{enumerate}
\renewcommand{\theenumi}{\roman{enumi}}
\renewcommand{\labelenumi}{(\theenumi)}
\item 假设 $p$ 保守. 对纤维 $E'_X$ 中的每个箭头 $u$,
$p(u)=\id_X$ 是同构, 所以 $u$ 在 $E'$ 中是同构,
从而在 $E'_X$ 中也是同构
(因为 $u$ 在 $E'$ 中的逆显然也在 $E'_X$ 中作逆).
因此 $E'_X$ 是群胚. 反过来, 假设 $E'_X$ 都是群胚,
设 $u'$ 是 $E'$ 中的一个箭头, 且 $p(u')$ 是同构; 我们证明 $u$ 是同构.
为此注意, 由于 $p$ 纤维化, 可将 $u':X'\rightarrow Y'$ 分解为复合
$X'\rightarrow u^\ast(Y')\rightarrow Y'$, 其中第一个箭头是 $X$-态射
(注意 $X=p(X')$, $u=p(u')$), 第二个是 $u$ 上的 Cartesian 态射.
第一个箭头是同构, 因为 $E'_X$ 是群胚;
第二个也是同构, 因为纤维化范畴的一个 \emph{Cartesian} 态射,
只要其投影是同构, 显然就是同构.

\item 假设 $p$ 忠实, 并设 $X'$, $Y'$ 是纤维范畴 $E_X$ 的两个对象.
从 $X'$ 到 $Y'$ 的任意两个箭头都位于 $E$ 的同一箭头 $\id_X$ 之上,
所以相等, 因而 $E_X$ 有序. 反过来,\pageoriginale 假设纤维范畴都有序,
我们证明 $p$ 忠实. 设 $u'$, $v':X'\rightrightarrows Y'$ 是 $E'$ 中的箭头,
位于 $E$ 的同一箭头 $u:X\rightarrow Y$ 之上.
它们便分解为 $X'\rightrightarrows u^\ast(Y)\rightarrow Y'$,
其中两个箭头 $X'\rightrightarrows u^\ast(Y')$ 是 $E'_X$ 中具有相同源和靶的箭头;
因此它们相等, 从而 $u'=v'$, \cqfd
\end{enumerate}

现在回到 \hyperref[I.prop6.5]{6.5} 的证明. 我们已证明
(i) $\Leftrightarrow$ (ii). 另一方面, (ii) $\Leftrightarrow$ (iv) 相当清楚:
事实上, 一方面, 范畴 $E_{/F}$ 在 $E$ 上纤维化,
它的纤维范畴是集合 $F(X)$ 所定义的离散范畴, 这直接由定义得出;
另一方面, 若 $p$ 如 (ii) 所述, 则由 SGA 1 VI 8 中的成串推论,
$E$ 上的纤维化范畴 $E'$ 与由函子 $E\rightarrow (\Cat)$
定义的 $E$ 上的分裂范畴 $E$-等价,
其中该函子由函子 $F:E\rightarrow (\Ens)$ 确定,
后者将每个 $X\in \ob E$ 对应到 $E'_X$ 的对象同构类的集合.
这个分裂范畴与范畴 $E_{/F}$ $E$-同构.
由于 (iv) $\Rightarrow$ (iii) 显然, 只须证明 (iii) $\Rightarrow$ (i).
显然, $p$ 忠实 (\resp 保守) 的充要条件是诱导函子
$E'_{/X'}\rightarrow E_{/p(X)}$ 也如此;
因此, 这些函子全忠实当然就足够了; 故只剩下证明 (iii) 蕴含 $p$ 纤维化.
但同样可见, 函子 $p$ 纤维化, 当且仅当诱导函子
$E'_{/X'}\rightarrow E_{/p(X')}$ 纤维化.
特别地, 当它们是在对象上满射的范畴等价时, 这一条件成立.


% END INCLUDED SGA4_English_current_buildable_source_20260906/source/01/drafts/SGA4_I_LINES_1947_2034_English.texfrag


% BEGIN INCLUDED SGA4_English_current_buildable_source_20260906/source/01/drafts/SGA4_I_LINES_2035_2353_English.texfrag
\section{生成子范畴与余生成子范畴}\etiquette[7]{I.sec7}\pageoriginale

\begin{enonce*}{定义 7.1}\etiquette[7.1]{I.def7.1}
设 $E$ 是一个范畴, $C$ 是 $E$ 的一个全子范畴.
若对 $E$ 的每个对象 $X$, $E$ 中以 $X$ 为靶、以 $X'\in \ob C$ 为源的箭头族
是满态射族 (\resp 严格满态射族) (\hyperref[I.condef1.3]{1.3}),
则称 $C$ 是 $E$ 的一个\emph{由严格满态射生成 (\resp 由满态射生成)} 的子范畴.
若对每个箭头 $u:Y\rightarrow X$
(\resp $E$ 的每个单态射, \resp $E$ 的每个严格单态射 $\hyperref[I.sec10.5]{(10.5)}$),
只要对每个 $X'\in \ob C$, 相应映射 $\Hom(X',Y)\rightarrow \Hom(X',X)$
是双射, $u$ 就是同构, 则称 $C$ 是 $E$ 的一个\emph{生成
(\resp 对单态射生成, \resp 对严格单态射生成)} 子范畴.
最后, 若由 $E$ 的对象族 $(X_i)$ 生成的 $E$ 的全子范畴 $C$
是由严格满态射生成 (\resp \ldots) 的,
则称这个对象族由严格满态射生成 (\resp \ldots).
\end{enonce*}

\setcounter{subsection}{1} \setcounter{subsubsection}{0}
\subsubsection{}\etiquette[7.1.1]{I.sec7.1.1}
注意, 用 $X'\in \ob C$ 所代表的函子族 $(h_{X'})_{X'\in \ob C}$
$$
h_{X'}:E\longrightarrow (\Ens)\; (X'\in \ob C)
$$
来表述, $C$ 生成 (\resp 对单态射生成, \resp 对严格单态射生成) 的条件,
就是说函子族 $(h_{X'})$ \emph{保守}
(\resp 对单态射保守, \resp 对严格单态射保守)
(\hyperref[I.def6.1]{6.1}). 由定义还直接可知,
$C$ 由满态射生成, 当\pageoriginale 且仅当族 $(h_{X'})_{X'\in \ob C}$
\emph{忠实} (\hyperref[I.def6.1]{6.1}).
下文 (\hyperref[I.prop7.2]{7.2} (i)) 还将对 $C$ 由严格满态射生成这一条件给出类似的解释.

\subsubsection{}\etiquette[7.1.2]{I.sec7.1.2}
与 \hyperref[I.def6.1]{6.1} 中引入的概念一样,
当 $E$ 有适当的正合性性质时, \hyperref[I.def7.1]{7.1} 中的概念尤其有用;
在这种情况下, 所引入的各个概念很有可能全部等价
(\hyperref[I.cor7.3]{7.3}).
因此, 几乎无须讨论 \hyperref[I.def7.1]{7.1} 中的这些概念哪一个最为重要;
在最重要的情形中, 这些概念重合,
因而“生成子范畴”这一术语可以不加区别地解释为
\hyperref[I.def7.1]{7.1} 中考虑的任一种性质
(例如第一种; 如我们将看到的 (\hyperref[I.prop7.2]{7.2} (ii)), 它是所有性质中最强的).

\subsubsection{}\etiquette[7.1.3]{I.sec7.1.3}
假设 $E$ 是一个 $𝒰$-范畴, 考虑典范函子
\begin{equation*}
\varphi:E\longrightarrow
  \widehat{C}=\Hhom(C, 𝒰\text{-}\Ens)\tag{7.1.3.1}
\end{equation*}\etiquette[7.1.3.1]{I.sec7.1.3.1}
它是函子 $E\rightarrow \widehat{E}\rightarrow \widehat{C}$ 的复合,
其中第一个是典范函子 (\hyperref[I.sec1.3.3]{1.3.3}),
第二个是限制到 $C$ 的函子.
注意, 说上述函子 $\varphi$ 保守 (\resp 忠实),
显然等价于说随着 $X'\in\ob C$ 变化的函子族
$h_{X'}:X\rightarrow \Hom(X',X)=\varphi(X)(X')$
是一个保守 (\resp 忠实) 族,
也就是 (\hyperref[I.sec7.1.2]{7.1.2}) 说 $C$ 生成 (\resp 由满态射生成).
类似地, $\varphi$ 对单态射\pageoriginale (\resp 严格单态射) 保守,
当且仅当 $h_{X'}(X'\in \ob C)$ 所成的族对单态射
(\resp 严格单态射) 保守,
\ie 当且仅当 $E$ 的子范畴 $C$ 对单态射 (\resp 严格单态射) 生成.

\begin{enonce*}{命题 7.2}\etiquette[7.2]{I.prop7.2}
设 $E$ 是一个 $𝒰$-范畴, $C$ 是一个全子范畴.
\begin{enumerate}
\renewcommand{\theenumi}{\roman{enumi}}
\renewcommand{\labelenumi}{(\theenumi)}
\item 下列条件等价:
\begin{enumerate}
\renewcommand{\theenumii}{\alph{enumii}}
\renewcommand{\labelenumii}{\theenumii)}
\item $C$ 是由严格满态射生成的子范畴.

\item 对每个 $X\in \ob E$, 用 $C_{/X}$ 表示 $E_{/X}$ 的全子范畴,
其对象为源 $X'\in \ob C$ 的箭头 $X'\rightarrow X$.
从包含函子 $j:C_{/X}\rightarrow E$ 到 $C_{/X}$ 上值为 $X$ 的常值函子的自然箭头,
使 $X$ 成为 $j$ 的归纳极限:
$$
X\xleftarrow{\sim} \lim_{C_{/X}}X'{\quoi}.
$$

\item (\hyperref[I.sec7.1.3.1]{7.1.3.1}) 中的典范函子 $\varphi$ 全忠实.
\end{enumerate}

\item \hyperref[I.def7.1]{7.1} 中的概念之间存在如下蕴含关系:
\[
\xymatrix@C=-.7cm@R=.45cm{
  &  \txt{{\rm 1)} $C$ 由严格满态射\\ 生成\\
  \emph{($\varphi$ 全忠实)}} \ar@{=>}[rd]&\\
 \txt{{\rm 2)} $C$ 由满态射\\ 生成\\
 \emph{($\varphi$ 忠实)}} &    &
 \txt{{\rm 3)} $C$ 生成 \phantom{aaa}\\
 \emph{($\varphi$ 保守)}} \ar@{=>}[d]\\
    &      &  \txt{{\rm 4)} $C$ 对单态射\\ 生成\\
  \emph{($\varphi$ 对单态射保守)}}\ar@{=>}[d]\\
    &      &  \txt{{\rm 5)} $C$ 对严格单态射\\ 生成\\
  \emph{($\varphi$ 对严格单态射保守).}}\\
}
\]

\item 此外\pageoriginale 还有下列带条件的蕴含关系:
\begin{enumerate}
\renewcommand{\theenumii}{\alph{enumii}}
\renewcommand{\labelenumii}{\theenumii)}
\item 若 $E$ 中的满态射族都是严格满态射族, 则 $2)\Rightarrow 1)$.
若 $E$ 中的单态射都是严格的, 则 $5)\Rightarrow 4)$.

\item 若 $E$ 中箭头对的核 (\resp 纤维积) 可表,
则 $3)\Rightarrow 2)$ (\resp $4)\Rightarrow 3)$).

\item 若 $E$ 中具有共同靶 $X$ 的每个态射族 $X_i\rightarrow X$
都可分解为严格满态射族 (\resp 满态射族) $X_i\rightarrow Y$,
后接一个单态射 (\resp 严格单态射) $Y\rightarrow X$,
则 $4)\Rightarrow 1)$ (\resp $5)\Rightarrow 2)$).
\end{enumerate}
\end{enumerate}
\end{enonce*}
立即记录如下推论.

\begin{enonce*}{推论 7.3}\etiquette[7.3]{I.cor7.3}
在以下两种情形中, $\hyperref[I.def6.1]{6.1}$ 中考虑的全部概念
(并在上面 {\rm(ii)} 的蕴含关系图中再次列出) 都等价:
\begin{enumerate}
\renewcommand{\theenumi}{\roman{enumi}}
\renewcommand{\labelenumi}{(\theenumi)}
\item $E$ 中箭头对的核和纤维积可表, 单态射都是严格的,
且满态射族都是严格满态射族.

\item $E$ 中具有共同靶 $X$ 的每个箭头族 $(X_i\rightarrow X)_{i\in I}$
都可分解为一个满态射族 $(X_i\rightarrow Y)$,
后接一个单态射 $Y\rightarrow X$,
$E$ 的每个单态射都严格, 且 $E$ 的每个满态射族都严格.
\end{enumerate}
\end{enonce*}

事实上,\pageoriginale 在情形 (i) 中, 由 b) 有 4) $\Rightarrow$ 3)
和 3) $\Rightarrow$ 2), 由 a) 有 5) $\Rightarrow$ 4)
和 2) $\Rightarrow$ 1). 在情形 (ii) 中, 由 c) 有蕴含关系
4) $\Rightarrow$ 1) 和 5) $\Rightarrow$ 2).
因此, 根据蕴含关系图 \hyperref[I.prop6.2]{6.2} (ii) 即得结论.

\medskip
\noindent \emph{\hyperref[I.prop7.2]{7.2} 的证明.}
\begin{enumerate}
\renewcommand{\theenumi}{\roman{enumi}}
\renewcommand{\labelenumi}{(\theenumi)}
\item 蕴含关系 b) $\Rightarrow$ a) 直接由定义得出.
我们证明 a) $\Rightarrow$ b). 因此在假设 a) 下,
须证明对所有 $X$, $Y\in \ob E$, 以 $X'\in \ob C_{/X}$ 为指标、
且对 $C_{/X}$ 中每个箭头 $f:X''\rightarrow X'$ 都满足
$u_{X'}f=u_{X''}$ 的每个箭头系
$$
u_{X'}:X'\longrightarrow Y
$$
都经过一个箭头 $X\rightarrow Y$ 分解 (由假设 a), 此箭头必唯一).
由假设 a), 只须验证对 $E_{/X}$ 的每个对象 $Z$ 和 $E_{/X}$ 中每对态射
$v':Z\rightarrow X'$, $v'':Z\rightarrow X''$,
其中 $X'$ 和 $X''$ 在 $C_{/X}$ 中, 有 $u_{X'}v'=u_{X''}v''$.
但由假设 a), 源满足 $X'''\in \ob C$ 的箭头族 $w:X'''\rightarrow Z$
是满态射族, 因此只须验证对每个这样的 $w$, 都有
$(u_{X'}v')w=(u_{X''}v'')w$.
这个等式也可写为 $u_{X'}(v'w)=u_{X''}(v''w)$,
而它直接由关于 $u$ 所成族的假设得出.

现在证明条件 b) 与 c) 等价. 为此注意, 对每个 $X\in \ob E$,
$\widehat{C}$ 的对象 $\varphi(X)$ 是典范函子
$\widehat{C}_{/\varphi(X)}\rightarrow \widehat{C}$ 在 $\widehat{C}$ 中的归纳极限
(\hyperref[I.prop3.4]{3.4});
但 $\widehat{C}_{/\varphi(X)}$ 典范同构于 $C_{/X}$,
所以在 $\widehat{C}$ 中有
$$
\varphi(X)=\varinjlim_{C_{/X}}X'{\quoi},
$$
其中\pageoriginale 将 $C$ 的对象 $X'$ 与它所代表的函子
$\in \widehat{C}$ 等同. 因此, 对 $E$ 的另一个对象 $Y$, 有典范同构
\begin{align*}
\Hom(\varphi(X),\varphi(Y))\simeq\varprojlim_{C_{/X}}
\Hom(X',\varphi(Y))&\simeq
\varprojlim_{C_{/X}}\varphi(Y)(X')\hyperref[I.prop1.4]{(1.4)}\\
&\simeq \varprojlim_{C_{/C}} \Hom(X'%%
,Y){\quoi}.\tag{$\ast$}
\end{align*}
于是, 通过 $(\ast)$ 两端之间的同构, 由 $\varphi$ 定义的映射
$$
\Hom(X,Y)\longrightarrow \Hom(\varphi(X),\varphi(Y))
$$
恰是由 \hyperref[I.prop7.2]{7.2} (i) b) 中考虑的、
以 $C_{/X}$ 为指标的箭头归纳系 $X'\rightarrow X$ 得出的映射
$$
\Hom(X,Y)\longrightarrow \varprojlim_{C_{/X}} \Hom(X',Y)
$$
因此, 固定 $X$ 后, 第一个映射对每个 $Y$ 都是双射,
当且仅当 $X$ 是包含函子 $j:C_{/X}\rightarrow E$ 的归纳极限;
这就证明了 b) 和 c) 等价.

\item 蕴含关系 1) $\Rightarrow$ 2) 和 3) $\Rightarrow$ 4)
$\Rightarrow$ 5) 由定义显然成立.
对于蕴含关系 1) $\Rightarrow$ 3), 根据 (i),
将性质 1) 解释为 $\varphi$ 全忠实, 再注意全忠实函子是保守的, 即可得出.
而我们已经指出 (\hyperref[I.sec7.1.1]{7.1.1}),
3) 就是说 $\varphi$ 保守.

\item 断言\pageoriginale a) 是重言式.
断言 b) 是将 \hyperref[I.prop6.2]{6.2} (i)
(\resp \hyperref[I.prop6.2]{6.2} (iv)) 应用于函子 $\varphi$ 所得的结论,
其中用到后者左正合这一事实. 最后证明 c).
对 $X\in \ob E$, 考虑满足 $X_i\in \ob C$ 的所有态射 $X_i\rightarrow X$
构成的族; 由关于 $E$ 的假设, 它分解为一个有效满态射族
(\resp 满态射族) $X_i\rightarrow Y$,
后接一个单态射 (\resp 严格单态射) $Y\rightarrow X$.
此时假设 4) (\resp 5)) 蕴含 $Y\rightarrow X$ 是同构,
所以所考虑的族是严格满态射族 (\resp 满态射族), \cqfd
\end{enumerate}

\begin{enonce*}{命题 7.4}\etiquette[7.4]{I.prop7.4}
设 $E$ 是一个范畴, $C$ 是 $E$ 的一个全子范畴,
$X$ 是 $E$ 的一个对象. 假设 $C$ 在 $E$ 中生成
(\resp $C$ 对单态射生成, 并且 $X$ 的两个子对象在 $X$ 上的纤维积在 $E$ 中可表).
则只要对每个 $T\in \ob C$, 已知 $\Hom(T,X')$ 的像所成的
$\Hom(T,X)$ 的子集, 就确定了 $X$ 的一个严格子对象
$\hyperref[I.sec10.11]{(10.11)}$ (\resp 子对象) $X'$.
因此, $X$ 的严格子对象集合 (\resp 子对象集合) 的基数
不超过 $\prod_{T\in \ob C} 2^{\card \Hom(T,X)}$;
若 $X\in \ob C$, 则不超过 $2^{\card%%
F\ell\,C}$.
\end{enonce*}

先证明括号内的断言. 设 $X'$ 和 $X''$ 是 $X$ 的两个子对象,
满足对每个 $T\in \ob C$, $\Hom(T,X')$ 和 $\Hom(T,X'')$
在 $\Hom(T,X)$ 中的像相等.
因此, 它们也等于 $\Hom(T,X''')$ 的像,
其中 $X'''$ 是 $X'$ 和 $X''$ 在 $X$ 上的纤维积.
由于 $C$ 对单态射生成, 可知单态射\pageoriginale
$X'''\rightarrow X'$ 和 $X'''\rightarrow X''$ 都是同构,
所以 $X'$ 和 $X''$ 相等, 两者均等于 $X$ 的子对象 $X'''$.

再证明不带括号的断言. 根据严格子对象这一概念的定义,
只须验证: 若对每个 $T\in\ob C$, 已知 $\Hom(T,X)$ 的子集 $\Hom(T,X')$,
就知道哪些箭头对 $\xymatrix{
X\ar@<2pt>[r]^{u,v}\ar@<-2pt>[r]& T }$ 满足 $ui=vi$,
其中 $i:X'\rightarrow X$ 是典范单射.
但由于 $C$ 生成, 关系 $ui=vi$ 等价于对每个 $f\in \Hom(T,X')$
都有 $(ui)f=(vi)f$,
\ie 对来自 $\Hom(T,X')$ 的每个 $g\in \Hom(T,X)$
(\ie 形如 if, 其中 $f\in \Hom(T,X')$), 都有 $ug=vg$,
这就证明了我们的断言.

\begin{enonce*}{推论 7.5}\etiquette[7.5]{I.cor7.5}
设 $E$ 是一个范畴, $C$ 是一个全生成子范畴, $X$ 是 $C$ 的一个对象.
则只要对每个 $T\in \ob C$, 已知由满足 $qu=qv$ 的对 $(u,v)$
构成的 $\Hom(T,X)^2$ 的子集, 就确定了 $X$ 的一个严格商
(\hyperref[I.sec10.8]{10.8}) $X'$, 其中 $q:X\rightarrow X'$ 是典范态射.
因此, $X$ 的严格商的集合的基数不超过
$\prod_{T\in \ob C} 2^{\card\Hom(T,X)^2}$.
\end{enonce*}

事实上, 根据定义, 只要对 $E$ 的每个对象 $Y$,
已知由满足 $qu=qv$ 的对 $(u,v)$ 组成的
$\Hom(Y,X)\times \Hom(Y,X)$ 的子集,
就确定了 $X$ 的一个严格商 $X'$, 其中 $q:X\rightarrow X'$ 是典范态射.
但由于 $C$ 生成, 关系 $qu=qv$ 等价于对每个源 $T$ 在 $C$ 中的态射
$f:T\rightarrow Y$, 都有 $(qu)f=(qv)f$.
这个关系也可写为 $q(uf)=q(vf)$, 这就证明了
\hyperref[I.cor7.5]{7.5} 的第一个断言. 第二个断言立即得出.

\setcounter{subsection}{5} \setcounter{subsubsection}{0}

% END INCLUDED SGA4_English_current_buildable_source_20260906/source/01/drafts/SGA4_I_LINES_2035_2353_English.texfrag


% BEGIN INCLUDED SGA4_English_current_buildable_source_20260906/source/01/drafts/SGA4_I_LINES_2354_2668_English.texfrag
\subsubsection{}\etiquette[7.5.1]{I.sec7.5.1}
我们\pageoriginale 可以推广 \hyperref[I.cor7.5]{7.5}:
对 $E$ 的一个对象族 $(X_i)_{i\in I}$,
引入\emph{该族在 $E$ 中的严格商}这一概念,
其含义是 $E$ 的态射组成的一个严格满态射族
$(p_i:X_i\rightarrow X')_{i\in I}$;
和通常的商一样, 这里约定, 若能找到一个同构
$v:X'\rightarrow X''$ (它必唯一),
使对每个 $i\in I$ 都有 $fp_i=q_i$,
就将两个这样的族 $(p_i:X_i\rightarrow X')$ 和
$(q_i:X_i\rightarrow X'')$ 等同.
采用这个术语后, \hyperref[I.cor7.5]{7.5} 的证明原封不动地适用,
给出如下变式:

\begin{enonce*}[remark]{变式 7.5.2}\etiquette[7.5.2]{I.var7.5.2}
设 $E$, $C$ 如 \hyperref[I.cor7.5]{7.5} 所述,
$(X_i)_{i\in I}$ 是 $E$ 的一个对象族.
则 $(X_i)_{i\in I}$ 在 $E$ 中的一个严格商 $X'$
(\hyperref[I.sec7.5.1]{7.5.1}) 由以下数据确定:
对每个 $T\in \ob C$ 和每对 $(i,j)\in I\times I$,
由满足 $p_iu=p_jv$ 的对 $(u,v)$ 组成的
$\Hom(T,X_i)\times \Hom(T,X_j)$ 的子集,
其中对每个 $i\in I$, $p_i:X_i\rightarrow X'$ 表示典范态射.
因此, $(X_i)_{i\in I}$ 在 $E$ 中的严格商的集合的基数不超过
$ \prod_{T\in \ob C} \prod_{i, j\in I}2^{\card \Hom(T,X_i)\times
  \card\Hom(T,X_j)}{\quoi}.
$
\end{enonce*}

\setcounter{subsubsection}{2}
\subsubsection{}\etiquette[7.5.3]{I.sec7.5.3}
显然, 若只假设 $C$ 对严格单态射生成, 类似的结论仍成立,
只要积 $X_i\times X_j$ 可表,
并且仅考虑族 $(X_{i})_{i\in I}$ 的\emph{有效}商,
\ie 满足纤维积 $\fprod{X_i}{X_{j}}{X'}$ 在 $E$ 中可表的严格商 $X'$
(若 $E$ 对纤维积稳定, 这就不构成限制).

\begin{enonce*}{命题 7.6}\etiquette[7.6]{I.prop7.6}
设\pageoriginale $E$ 是一个范畴,
$C$ 是 $E$ 的一个由满态射生成的全子范畴
(\hyperref[I.def7.1]{7.1}),
$\Pi_\circ=$ 是一个无限基数 $\geq \card  F\ell  C$,
$\Pi\geq \Pi_\circ$ 是一个基数,
$(u_i:X_i\rightarrow X)_{i\in I}$ 是 $E$ 中一个普遍满态射族
(\hyperref[I.sec10.3]{10.3}),
由可作基变换的态射 (\hyperref[I.sec10.7]{10.7}) 组成,
源满足 $X_i\in\ob C$, 且 $\card I\leq \Pi$.
则 $\card F\ell  C_{/X}\leq \Pi^{\Pi_\circ}$.
\end{enonce*}

设 $T\in \ob C$; 只须证明
\begin{equation*}
\card\Hom(T,X)\leq
  \Pi^{\Pi_\circ}{\quoi},\tag{7.6.1}
\end{equation*}\MakeLinkTarget*{sga4.exp.I.equation.7.6.1}\label{I.eq7.6.1}
因为由此可依次得到
$$
\card\ob C_{/X}\leq(\Pi^{\Pi_\circ})\times\card \ob C_{/X}\leq
(\Pi^{\Pi_\circ})\times \Pi_\circ=\Pi^{\Pi_\circ}{\quoi},
$$
最后得到 $\card F\ell C_{/X}\leq ((\Pi^{\Pi_\circ})^2)\times
\Pi_\circ=\Pi^{\Pi_\circ}$,
因为对 $C_{/X}$ 的两个对象 $S$, $T$, 有
$\card \Hom_{C_{/X}}(S,T)\leq
\car\Hom_C(S,T)\leq\Pi_\circ$.
为证明 \ref{I.eq7.6.1}, 先注意如下引理:

\begin{enonce*}{引理 7.6.2}\etiquette[7.6.2]{I.lem7.6.2}
设 $J$ 是一个满足 $\card J=\Pi_\circ$ 的集合.
对 $C$ 的每个对象 $T$ 和每个态射 $f:T\rightarrow X$,
存在一个满态射族 $(v_j:S_j\rightarrow T)_{j\in J}$,
源满足 $S_j\in\ob C$,
并且对每个 $j\in J$, 存在 $i(j)\in I$ 和
$g_j:S_j\rightarrow X_{i(j)}$, 使 $u_{i(j)}g_j=fv_j$.
\end{enonce*}

事实上, 由假设, $\fprod{X_i}{T}{X}\rightarrow$ 所成的族是满态射族;
另一方面, 由于 $C$ 由严格满态射生成, 对每个 $i$,
源满足 $S\in \ob C$ 的箭头 $S\rightarrow \fprod{X_i}{T}{X}$
所成的族是满态射族.
因此, 由传递性, 源在 $\ob C$ 中、并经过某个 $\fprod{X_i}{T}{X}$
分解的箭头 $v:S\rightarrow T$ 所成的族是满态射族.
这些恰好是源在 $C$ 中且满足存在 $i\in I$ 和 $g:S\rightarrow X_i$
使 $u_ig=fv$ 的箭头 $v:S\rightarrow T$.
由于\pageoriginale 这些箭头 $v:S\rightarrow T$ 组成的集合包含在 $F\ell$ $C$ 中,
从而基数 $\leq \Pi_\circ$, 所以它可以用指标集 $J$ 编号, 引理得证.

现在注意, 知道了 $fv_j$ 就确定了态射 $f:T\rightarrow X$,
而知道了 $g_j$ 又确定了它们.
因此, $\card\Hom(T,X)$ 不超过族 $(i(j),S_j,v_j,g_j)_{j\in J}$
组成的集合的基数;
由于映射 $j\rightarrow I$ 组成的集合的基数 $\leq \Pi^{\Pi_\circ}$,
并且固定这样的一个映射 $j\mapsto i(j)$ 后,
相应的族 $S_j$, $f_j$, $v_j$ 组成的集合的基数,
不超过从 $J$ 到 $F\ell  C\times F\ell  C$ 的映射集的基数,
而后者 $\leq \Pi^{\Pi_\circ}_\circ$,
因为 $\card(F\ell C\times F\ell C)=\Pi_\circ^2=\Pi_\circ$,
所以 \hyperref[I.eq7.6.1]{(7.6.1)} 的左端不超过
$\Pi^{\Pi_\circ}\times \Pi_\circ^{\Pi_\circ}=\Pi^{\Pi_\circ}$,
这就完成了 \hyperref[I.prop7.6]{7.6} 的证明.

\begin{enonce*}{命题 7.7}\etiquette[7.7]{I.prop7.7}
设 $E$ 是一个纤维积可表的范畴,
$(X_\alpha)_{\alpha\in A}$ 是 $E$ 的一个生成对象族,
$(\varphi_i)_{i\in I}$ 是一个与纤维积交换的函子族
$\varphi_i:E\rightarrow E_i$.
族 $(\varphi_i)$ 保守 (\hyperref[I.prop5.1]{5.1}),
当且仅当对每个 $\alpha \in A$ 和 $X_\alpha$ 的每个不同于 $X_\alpha$ 的子对象 $X'$,
存在 $i\in I$ 使 $\varphi_i(X')\rightarrow \varphi_i(X_\alpha)$ 不是同构.
此时, 若 $E$ 是一个 $𝒰$-范畴且 $A$ 是 $𝒰$-小的,
则存在 $I$ 的一个 $𝒰$-小子集 $J$,
使 $(\varphi_j)_{j\in J}$ 已经是一个保守函子族.
\end{enonce*}

所述保守性条件的必要性显然, 我们证明其充分性.
由 \hyperref[I.def6.1]{6.1} (iv),
只须证明 $(\varphi_i)$ 对单态射保守.
设 $u:Y'\rightarrow Y$\pageoriginale 是 $E$ 中一个非同构的单态射;
须证明存在 $i\in I$ 使 $\varphi_i(u)$ 不是同构.
由关于 $(X_\alpha)$ 的假设,
存在 $\alpha$ 和一个不经过 $Y'$ 分解的态射 $X_\alpha\rightarrow Y$;
换言之, $Y'$ 的逆像 $X'$ 是 $X_\alpha$ 的一个不同于 $X_\alpha$ 的子对象.
由假设, 存在 $i\in I$ 使 $\varphi_i(X')\rightarrow \varphi(X)$ 不是同构.
由于 $\varphi_i(X')\simeq
\fprod{\varphi(X_\alpha)}{\varphi_i(Y')}{\varphi_{i(Y)}}$,
确可推出 $\varphi_i(Y')\rightarrow \varphi_i(Y)$ 不是同构.

根据 \hyperref[I.prop7.4]{7.4},
\hyperref[I.prop7.7]{7.7} 的第二个断言直接由第一个断言得出.

\begin{enonce*}{推论 7.7.1}\etiquette[7.7.1]{I.cor7.7.1}
设 $E$ 是一个 $𝒰$-范畴, 其中纤维积可表,
并且它有一个 $𝒰$-小的生成对象族.
则 $E$ 的每个生成族 $(Y_i)_{i\in I}$
都包含一个 $𝒰$-小的生成子族 $(Y_j)_{j\in J}(\card J\in 𝒰)$.
\end{enonce*}

事实上, 只须将 \hyperref[I.prop7.7]{7.7} 应用于
$E$ 的一个 $𝒰$-小生成族 $(X_\alpha)_{\alpha\in A}$
以及由 $Y_i$ 代表的函子族 $\varphi_i$ $(i\in I)$.

\begin{enonce*}{命题 7.8}\etiquette[7.8]{I.prop7.8}
设 $E$ 是一个范畴, $C$ 是一个由严格满态射生成的全子范畴,
$D$ 是一个范畴,
$\Hhom'(E,D)$ 是 $\Hhom(E,D)$ 的全子范畴,
由与 $C_{/X}$ 型归纳极限交换的函子组成, 其中 $X$ 是 $E$ 的任意对象.
则函子 $F\rightarrow F| C$ 诱导一个全忠实函子
$$
\Hhom'(E,D)\longrightarrow \Hhom(C,D){\quoi}.
$$
因此,\pageoriginale 若 $\Hhom(C,D)$ 是一个 $𝒰$-范畴
$(\hyperref[I.def1.1]{1.1}.)$,
例如当 $C$ 是 $𝒰$-小的且 $D$ 是 $𝒰$-范畴时
{\rm(\hyperref[I.sec1.1.1]{1.1.1}. b))},
则 $\Hhom'(E,D)$ 也是一个 $𝒰$-范畴.
\end{enonce*}

设 $F$, $G:E\rightarrow D$ 是左端范畴中的两个函子,
$u:F\rightarrow G$ 是一个自然变换.
若在 $E$ 中 $X=\varinjlim X_i$,
且 $F$ 和 $G$ 与所考虑的归纳极限交换,
则 $u(X):F(X)\rightarrow G(X)$ 可视为态射
$u(X_i):F(X_i)\rightarrow G(X_i)$ 的归纳极限,
所以知道了 $u(X_i)$ 就确定了它.
根据 \hyperref[I.prop7.2]{7.2} (i) 中的蕴含关系 a) $\Rightarrow$ b),
这表明 \hyperref[I.prop7.8]{7.8} 中所考虑的函子忠实.
反过来, 设 $v:F|\sisi{D}{C}\rightarrow G|\sisi{D}{C}$ 是一个自然变换;
我们证明它来自一个自然变换 $u:F\rightarrow G$.
对每个 $X\in \ob E$, 将
$$
u(X):F(X)=\varinjlim_{C_{/X}}F(X_i)\longrightarrow
G(X)=\varinjlim_{C_{/X}}G(X_i)
$$
定义为 $v(X_i)$ 的归纳极限.
显然, 如此得到的映射关于 $X$ 具有函子性,
因此给出一个自然变换 $u:F\rightarrow G$,
\sisi{中}{并且}最后, $u$ 从 $F|C\sisi{'}{}$ 到 $G|C\sisi{'}{}$
诱导的态射就是 $v$, 证明完成.


\setcounter{subsection}{8}
\subsection{余生成族与余生成子范畴}\etiquette[7.9]{I.sec7.9}

设 $E$ 是一个范畴, $C$ 是 $E$ 的一个全子范畴.
若 $E^\circ$ 的全子范畴 $C^\circ$ 由严格满态射生成 (\resp\ldotsn),
则称 $C$ 由严格单态射余生成 (\resp 由单态射余生成,
\resp 余生成, \resp 对满态射余生成, \resp 对严格满态射余生成).
对象族的术语与此类似.
因此本节关于\pageoriginale 生成子范畴概念及其变式
(\hyperref[I.def7.1]{7.1}) 的全部结果,
都给出这些对偶概念的相应结果; 具体表述留待读者自行写出.

\begin{enonce*}{命题 7.10}\etiquette[7.10]{I.prop7.10}
设 $E$ 是一个范畴, $C$ 是一个全生成子范畴
$\hyperref[I.def7.1]{(7.1)}$, $D$ 是 $E$ 的一个全子范畴.
子范畴 $D$ 余生成 $\hyperref[I.sec7.9]{(7.9)}$,
当且仅当对 $E$ 中每对源满足 $T\in \ob C$ 且 $u\neq v$ 的平行箭头
$\xymatrix{ T\ar@<2pt>[r]^{u,v}\ar@<-2pt>[r]& X }$,
存在一个靶满足 $I\in \ob D$ 的箭头 $w:X\rightarrow I$, 使 $wu\neq wv$.
\end{enonce*}

根据定义, 说 $D$ 余生成, 就是说对 $E$ 中满足 $u\neq v$ 的每对平行箭头
$\xymatrix{ Y\ar@<2pt>[r]^{u,v}\ar@<-2pt>[r]& X
}$, 存在一个箭头 $w:X\rightarrow I$,
其中 $I\in\ob D$, 使 $wu\neq wv$.
因此 \hyperref[I.prop7.10]{7.10} 仅仅是说,
检验这个性质时只须考虑 $Y\in \ob C$ 的情形.
现在, 由于 $C$ 生成, 假设 $u\neq v$ 蕴含存在
$f:T\rightarrow Y$ 使 $uf\neq vf$;
于是由假设, 存在靶满足 $I\in\ob D$ 的 $w:X\rightarrow I$,
使 $w(uf)\neq w(vf)$, 从而 $wu=wv$, \cqfd

\begin{enonce*}{推论 7.11}\etiquette[7.11]{I.cor7.11}
沿用 \hyperref[I.prop7.10]{7.10} 的记号,
假设 $D$ 的对象 $I$ 都是 $E$ 的内射对象,
\ie 对 $E$ 中每个单态射 $X\rightarrow Y$,
相应映射 $\Hom(Y,I)\rightarrow \Hom(X,I)$ 都是满射.
还假设 $E$ 中每对平行箭头 $\xymatrix{
T\ar@<2pt>[r]^{u,v}\ar@<-2pt>[r]& X }$
都可分解为一对满态射 (\resp 有效满态射) 平行箭头
$\xymatrix{ Y\ar@<2pt>[r]^{u',v'}\ar@<-2pt>[r]& X }$,
后接一个单态射 $i:X'\rightarrow X$.
则在 \hyperref[I.prop7.10]{7.10} 中 $D$ 余生成的判据中,
可以仅考虑满态射 (\resp 有效满态射) 对 $(u,v)$.
\end{enonce*}

我们\pageoriginale 得到:

\begin{enonce*}{推论 7.12}\etiquette[7.12]{I.cor7.12}
设 $E$ 是一个 $𝒰$-范畴, 它有一个小生成子范畴 $C$,
并且 $E$ 的每个对象都能作为某个单态射的源, 该单态射的靶是 $E$ 的内射对象.
还假设对每个 $T\in\ob C$, $E$ 中的余积 $T\amalg T$ 可表,
且以 $T\amalg T$ 为源的每个态射都可分解为有效满态射后接单态射.
则 $E$ 有一个小的全余生成子范畴 $D$.
更确切地, 可以取 $D$ 使 $\card \ob D\leq
\prod_{T,T'\in \ob C}2^{\card(\Hom(T',T\coprod T))^2}$.
\end{enonce*}

事实上, 由 \hyperref[I.cor7.11]{7.11},
只须对每个 $T\in \ob C$ 以及 $T\amalg T$ 的每个有效商 $X$,
选取 $X$ 到 $E$ 的一个内射对象 $I$ 的嵌入,
再取 $D$ 为这些 $I$ 张成的 $E$ 的全子范畴.
于是由 \hyperref[I.cor7.5]{7.5} 得到结论.

\begin{enonce*}[remark]{例 7.13}\etiquette[7.13]{I.exem7.13}
因此, 要由小生成子范畴构造小余生成子范畴,
就须寻找 $𝒰$-范畴 $E$ 有“足够多内射对象”的条件,
\ie 每个对象都能 (通过单态射) 嵌入某个内射对象.
众所周知 [Tohoku], 若一个 Abel $𝒰$-范畴有正合的 (小) 滤过归纳极限
(\loccit 中的公理 $AB$ 5), 并且有一个小生成族, 则它满足这个条件.
此外, \loccit 中的构造并非本质上局限于 Abel 范畴,
也适用于拓扑空间 $X\in 𝒰$ 上的 $𝒰$-集合层范畴.
我们不在此陈述使这个构造有效的正合性性质,
只指出, 在\pageoriginale $X$ 上的集合层范畴这一特殊情形中,
可以如下直接归结到 \loccit 的情形.
注意, 若 $\mathcal{O}_X$ 是 $X$ 上的一个环,
则 $\mathcal{O}_X$-模范畴的每个内射对象 $I$,
作为集合层范畴的对象也是内射的.
事实上, 若 $F$ 是一个集合层,
$\mathcal{O}_X[F]$ 是“由 $F$ 生成的自由 $\mathcal{O}_X$-模”,
则根据定义, 有一个集合层同态
\begin{equation*}
F\longrightarrow \mathcal{O}_X[F]\tag{$\ast$}
\end{equation*}
它给出一个关于 $F$ 具有函子性的同构
$$
\Hom_{\mathcal{O}_X}(\mathcal{O}_X[F], I)\xrightarrow{\sim} \Hom
(F,I){\quoi}.
$$
由于函子 $F\mapsto \mathcal{O}_X[F]$ 显然把单态射映为单态射,
立即可知, 若 $I$ 是内射 $\mathcal{O}_X$-模,
它也是一个内射集合层.
由此可知, 若态射 $(\ast)$ 是单态射
(例如取 $\mathcal{O}_X$ 为值是某个环 $A\neq 0$ 的常值环时就是如此),
则 $\mathcal{O}_X[F^\ast]$ 到某个内射 $\mathcal{O}_X$-模 $I$ 的嵌入,
给出 $F$ 到 $I$ 的底层内射集合层的一个嵌入.
\end{enonce*}

同样的论证原封不动地适用于 $𝒰$-景上的 $𝒰$-集合层范畴,
这一范畴将在下一讲中引入.

小余生成子范畴的存在性之所以重要,
主要在于下文的可表性判据 \hyperref[I.cor8.12.7]{8.12.7}.

\setcounter{section}{7}

% END INCLUDED SGA4_English_current_buildable_source_20260906/source/01/drafts/SGA4_I_LINES_2354_2668_English.texfrag


% BEGIN INCLUDED SGA4_English_current_buildable_source_20260906/source/01/drafts/SGA4_I_LINES_2669_2936_English.texfrag
